\section{Steering: making the scaffold axis causal}
\label{app:steer}

\paragraph{Methods.}
\Cref{app:geometry} shows the \emph{scaffold} direction is the most separable axis
of the answer representation, but that is correlational. To test it causally we
turn the geometry into an additive intervention. From the geometry residuals
(\Cref{app:geometry}) we pair each item's scaffold and no-scaffold activation by
question id and take the per-layer \emph{diff-of-means}, $v^{(l)} =
\overline{h^{(l)}_{\text{scaffold}}} - \overline{h^{(l)}_{\text{none}}}$, as a
steering vector at the primary layer $l{=}14$ (depth $0.50$). We split the
$231$ pairs by question id into a seeded $70/30$ train/test split (no item
overlap) and \emph{fit nothing} beyond the train-mean direction; all numbers below
are on the held-out \textbf{test} items. At decode we add $\alpha\,v^{(14)}$ to the
\texttt{resid\_post} stream of every token and re-read the teacher-forced
three-option abstention on the $n{=}24$ ambiguous test items, sweeping
$\alpha \in [-2, 4]$. A positive $\alpha$ should \emph{raise} abstention (push
toward the scaffold); a negative $\alpha$ should suppress it.

\paragraph{Results.}
\begin{itemize}[leftmargin=*,topsep=2pt,itemsep=1pt]
  \item Adding the scaffold axis \emph{causally raises abstention} on held-out
    ambiguous items: test abstention climbs from $0.54$ at $\alpha{=}0$ to $0.71$
    at $\alpha{=}{+}1$ and a peak $0.75$ at $\alpha{=}{+}2$ (unknown-mass
    $0.49 \to 0.62$).
  \item The negative direction \emph{suppresses} it: $\alpha{=}{-}2$ drops
    abstention to $0.17$ (unknown-mass $0.18$), confirming a signed, monotone axis
    rather than a generic noise perturbation.
  \item The effect is bounded, not a runaway: over-steering at $\alpha{=}{+}4$
    collapses back to the $\alpha{=}0$ baseline ($0.54$), so the axis recovers part
    of---but does not saturate to---the scaffold ceiling.
  \item Reference ceiling: the actual scaffold prompt reaches $1.00$ abstention
    ($0.99$ unknown-mass) on the same $n{=}24$ test items, so a single mid-layer
    diff-of-means vector recovers a substantial fraction of the scaffold's
    behavioral effect with no prompt change.
\end{itemize}

\begin{figure}[htbp]\centering
  \begin{subfigure}{0.49\linewidth}\centering
    \plotsub{../out/sesgo/steer/figures/abstention_vs_alpha.png}{abstention_vs_alpha.png}{\linewidth}
    \caption{Held-out test abstention vs.\ steering strength $\alpha$: rises from
      $0.54$ ($\alpha{=}0$) to a peak $0.75$ ($\alpha{=}{+}2$), suppressed to
      $0.17$ at $\alpha{=}{-}2$.}
    \label{fig:steer-abstain}
  \end{subfigure}\hfill
  \begin{subfigure}{0.49\linewidth}\centering
    \plotsub{../out/sesgo/steer/figures/unknown_prob_vs_alpha.png}{unknown_prob_vs_alpha.png}{\linewidth}
    \caption{Mean three-option unknown-mass vs.\ $\alpha$ on the same test items
      ($0.49 \to 0.62$ peak; scaffold reference $0.99$).}
    \label{fig:steer-unknown}
  \end{subfigure}
  \caption{\textbf{Steering} on \texttt{Qwen3-0.6B} (layer $14$ diff-of-means,
    $231$ pairs $\to$ $70/30$ split, $n{=}24$ held-out ambiguous test items):
    adding $+\alpha v$ on the scaffold axis causally raises abstention; $-\alpha$
    suppresses it. This makes the \Cref{app:geometry} scaffold axis causal.}
  \label{fig:steer}
\end{figure}

%%%%%%%%%%%%%%%%%%%%%%%%%%%%%%%%%%%%%%%%%%%%%%%%%%%%%%%%%%%%%%%%%%%%%%%%%%
